This work had investigated the possibility to implement an inline edge coverage instrumentation for binary only 
grey-box fuzzing that uses GRAPE binary rewriter framework to add instructions to the binaries.
The instrumentation added is compatible with AFL and AFL++ and require no access to the source code.
The following chapter summarize the work discussed in the previous chapters, asseses to which of the goals 
in \cref{sec:goals} were achieved and close with some final remarks.
The limitations of the vurrent implementation have been already discussed in \cref{chap:limitations} and so are not 
repeated here.

\section{Summary of the work}\label{sec:summary}

    Coverage-guided grey-box fuzzing use the coverage feedback obtained through instructions instrumented inside the binary. 
    Dynamic approached such as AFL++ in QEMU or FRIDA mode, DynamoRIO \dots insert the coverage tracing instruction 
    dynamically within each \gls{BB}. This approach is easy to deploy but costly in terms of overheads.
    Static approaches, such the one implemented in this work, statically rewrites the binary introducing 
    the coverage tracing instruction, making the instrumented binary quite fast in term of overheads but with 
    difficulty in maintain correctness.

    The presented approach is based on the injection of a shared object (\texttt{libforkserver.so}) that implement 
    the communication protocol with the fuzzer, implement the forkserver logic and attach the shared memory to the process.
    In order to inject the shared object a \texttt{DT\_NEEDED} dependency is added to the binary, alongside with 
    some relocation entries and some new sections, used as a support to the instrumentation. 
    This has been implemented in the \texttt{instrument.py} file, togheter with the Instrument point search and the instrumntation.
    The binary modification and the actual instrumentation are made available thanks to two libraries.
    The binary modification (add dynamic dependency, relocations and sections) is done with LIEF python library.
    The actual instrumentation is made with the GRAPE binary rewriting framework.

\section{Assessment of the goals}\label{sec:goal_assesment}

    \begin{center}
        \textbf{Goal I - behavior preservation}
    \end{center} 
    The Experiment I was performed on five stripped binaries taken from the ZAFL dataset.
    Among these, only two of them (\texttt{bsdtar} and \texttt{sfconvert}) preserve their behavior on all the tested 
    input (34 out of 34 for \texttt{bsdtar} and 52 out of 52 for \texttt{sfconvert}).
    \texttt{readelf} was instrumented but fails to maintain the behavior. The cause has been found to be the \gls{ELF} 
    layout modification made automatically by LIEF that change the address in which some segments are loaded without 
    changin the instructions that reference these segments.
    This cause a segmentation fault each time an instruction try to access the old memory address.
    \texttt{readelf} pass 4 out of 54. The result is given by the fact that also the original binary terminate with a non-zero 
    exit code and the behavior of the original and the instrumneted is erroneously considered to be the same.
    Moreove, Experiment II reveal one additional failure on \texttt{who}, caused by an error on the instrumentation itself 
    that still has to be investigated.
    For all these reason, Goal I is considered only partially achieved. Infact the robustness of the pipeline presented 
    in not sufficiently strong to address any binary in the general case

    \begin{center}
        \textbf{Goal II - compatibility with AFL}
    \end{center} 
    Experiment II shown that forkserver injected into the binary can correctly communicate with AFL++ over 24-hours 
    fuzzing campaign.
    Both the forkserver handshake, the shared memory update and the collection of the exit status of the target 
    works correctly in each repetition  on \texttt{base64} and \texttt{uniq}.

    \begin{center}
        \textbf{Goal III -  overhead with respect to other binary-only approaches.}
    \end{center} 
    On the two correctly instrumented LAVA-M binaries (\texttt{base64} and \texttt{uniq}), the version instrumneted 
    with the tool presented in this thesis sustain on average about 1760 and 1670 executions per second respectively, 
    which is between 2.9 and 4.4 times the throughput of AFL++ in QEMU mode (which has about 400 and 570 executions per second 
    rispectively).
    With respect to ZAFL instrumnetation method, the throughput is of the same order and the difference is smaller than the 
    variation observed between repetitions of the same configuration.
    So neither an improvement nor a deterioration of the perfomrance could be claimed.
    Two caveat must be kept in mind while these results are read : 
    first of all the number of edges that, in the case of the tool discussed in this 
    work is considerably smaller with respect to AFL++ in QEMU mode and ZAFL.
    I had suppose that this is due to the choice to not instrument the calls, reducing the number of edges of the \gls{CFG} 
    that are actually been instrumented.
    The second is that the evaluation done rely on small number of binaries and does not include a baseline 
    for compile-time instrumentation, so the overheads could not be quantified in absolute terms.
    Therefore Goal III is achieved with respect to Dynamic Binary Instrumentation. 
    For the competitivity with other Static Binary Instrumentation methods the Goal III remains partially achieved.
    Infact the achieved performance are comparable with ZAFL but the fact that this last tool preserve the behavior 
    in more casesmake the developed tool not competitive for the moment with the state-of-the-art.

\section{Final remarks}

    In the end I have showed that an inline, stack-free edge coverage instrumentation build on GRAPE is feasible 
    and the obtained executions per second in a fuzzing campaign are well above \gls{DBI} and in average comparable 
    with the throughput of other \gls{SBI}.
    The experiments shows that the approach is currently weak. 
    Infact, although most of the failures does not originate from the instrumentation logic itself but from the 
    infrastructure around it, the modification of a binary adding relocations, sections and dependencies without 
    modify \gls{ELF} layout remain a problem with the proposed approach.
    This could be solved scanning the instructions searching the one that reference the old sections 
    and update their offset, in order to reference the same location of the new section.

    For all these reason, the developed tool should be consider a working prototype and a proof of concept 
    of the approach, rather than a general porpouse rewriter.
    Its performance results are promising but preliminary given the small number of binary on which has been tested.
